\documentclass[12pt]{article}
\usepackage{bbm}
\usepackage{fancyhdr}
\usepackage{titling}
\usepackage{enumitem}
\usepackage{amsmath}
\usepackage{dsfont}
\usepackage{amsfonts}
\usepackage{amssymb}
\usepackage{amsthm}
\usepackage{geometry}
\usepackage{tikz}
\usepackage{algorithm}
\usepackage{algpseudocode}
\geometry{verbose,tmargin=1in,bmargin=1in,lmargin=1in,rmargin=1in}
\renewcommand{\vec}[1]{\mathbf{#1}}
\newcommand{\dec}{\text{dec}}
\newcommand{\divi}{\mathbin{\textnormal{div}}}
\title{Homework 1}
\author{Satvik Prasad}
\date{}
\newtheorem{lemma}{Lemma}[section]
\fancypagestyle{plain}{%
  \fancyhf{}
  \fancyhead[L]{\theauthor}
  \fancyhead[R]{\thedate}
}
\pagestyle{plain}
\makeatletter
\def\@maketitle{%
  \newpage
  \null
  \vskip 1em%
  \begin{center}%
    \let \footnote \thanks
    {\LARGE \@title \par}%
    \vskip 1em%
  \end{center}%
  \par
\vskip 1em}
\makeatother
\usepackage{cmbright}

\newcommand{\mat}[1]{\mathbf{#1}}

% ---------- Probability theory ----------
% Core operators (optional argument-free; use as \E{X}, \Var{X}, \Cov{X}{Y})
\DeclareMathOperator{\EOp}{\mathbb{E}}
\DeclareMathOperator{\VarOp}{Var}
\DeclareMathOperator{\CovOp}{Cov}
\DeclareMathOperator{\CorrOp}{Corr}

\renewcommand{\P}[1]{\mathbb{P}\left(#1\right)}        % \P{A}
\newcommand{\E}[1]{\EOp\left[#1\right]}                % \E{X}
\newcommand{\Var}[1]{\VarOp\left(#1\right)}            % \Var{X}
\newcommand{\Cov}[2]{\CovOp\left(#1,\,#2\right)}       % \Cov{X}{Y}
\newcommand{\Corr}[2]{\CorrOp\left(#1,\,#2\right)}     % \Corr{X}{Y}
\newcommand{\SD}[1]{\operatorname{SD}\left(#1\right)}  % \SD{X}

% Subscripted versions, e.g. \Esub{\theta}{X}, \Psub{\theta}{A}
\newcommand{\Esub}[2]{\EOp_{#1}\left[#2\right]}
\newcommand{\Psub}[2]{\mathbb{P}_{#1}\left(#2\right)}

% Conditioning, indicators, independence
\newcommand{\given}{\,\middle|\,}                      % \E{X \given Y}
\newcommand{\ind}[1]{\mathbbm{1}\left\{#1\right\}}     % \ind{X > 0}
\newcommand{\indep}{\perp\!\!\!\perp}                  % X \indep Y
\newcommand{\iid}{\overset{\text{iid}}{\sim}}          % X_i \iid \Normal(0,1)

% Convergence
\newcommand{\convp}{\xrightarrow{\;p\;}}               % in probability
\newcommand{\convd}{\xrightarrow{\;d\;}}               % in distribution
\newcommand{\convas}{\xrightarrow{\text{a.s.}}}        % almost surely
\newcommand{\convLp}[1]{\xrightarrow{L^{#1}}}          % in L^p

% Probability spaces
\newcommand{\F}{\mathcal{F}}                           % sigma-algebra
\newcommand{\B}{\mathcal{B}}                           % Borel sigma-algebra
\newcommand{\pspace}{(\Omega, \F, \mathbb{P})}

% Common distributions
\newcommand{\Normal}{\mathcal{N}}
\newcommand{\Bern}{\operatorname{Bernoulli}}
\newcommand{\Bin}{\operatorname{Binomial}}
\newcommand{\Pois}{\operatorname{Poisson}}
\newcommand{\Unif}{\operatorname{Uniform}}
\newcommand{\Exp}{\operatorname{Exp}}
\newcommand{\Geom}{\operatorname{Geometric}}
\newcommand{\Gam}{\operatorname{Gamma}}
\newcommand{\Beta}{\operatorname{Beta}}

% Number sets
\newcommand{\R}{\mathbb{R}}
\newcommand{\N}{\mathbb{N}}
\newcommand{\Z}{\mathbb{Z}}


\begin{document}
\begin{enumerate}
  \item The pseudocode of such an algorithm may be described as
    \begin{algorithm}
      \caption{D\&C Max}
      \begin{algorithmic}
        \Procedure{Max}{$A[1 \ldots n]$}
        \If{$n = 1$}
        \State \Return $A[1]$
        \EndIf

        \State $L \gets \text{Max}(A[1\ldots n/2])$
        \State $R \gets \text{Max}(A[n/2 + 1 \ldots n])$

        \If{$L \geq R$}
        \State \Return $L$
        \EndIf

        \State \Return $R$
        \EndProcedure
      \end{algorithmic}
    \end{algorithm}

    We prove correctness inductively. When $n = 1,$ $A[1]$ is
    trivially the maximum element, meaning this algorithm is correct
    for input size $1.$

    Assume this algorithm is true for all input sizes $1 \leq k \leq
    n - 1.$ Then, for input size $n,$ the maximum either lives in
    $A[1\ldots n/2]$ or $A[n/2 + 1 \ldots n].$ Let the maximum
    element be $m = A[i].$ By our inductive hypothesis on $k = n/2,$
    $\text{Max}(A[1\ldots n/2]) = \max_{1 \leq j \leq n/2} A[j]$ and
    $\text{Max}(A[n/2+1\ldots n]) = \max_{n/2 < j \leq n} A[j]$
    \begin{itemize}
      \item If $1 \leq i \leq n/2,$ then $m = \max_{1 \leq j \leq
        n/2}A[j] \geq \max_{n/2 <j \leq n}A[j]$, since otherwise $m$
        cannot be the maximum. In this case, we have $L \geq R,$ and
        thus we return $L = m,$ meaning our algorithm is correct for size $n.$

      \item Alternatively, if $n/2 < i \leq n,$ then similarly $L =
        \max_{1 \leq j \leq n/2}A[j] \leq \max_{n/2 <j \leq n}A[j] =
        R = m.$ In this case, if $L = R = m$ then we take the $L \geq
        R$ branch and return $m.$ Otherwise, if $L < R,$ then we skip
        the branch and return $R = m.$ In either case, we return $m,$
        meaning our algorithm is correct for size $n$.
    \end{itemize}

    $\therefore$ by strong induction, this algorithm is correct for
    all input sizes $n \geq 1.$

    The recurrence of this algorithm satisfies $T(n) = 2T(n/2) +
    \mathcal O(1),$ and $T(1) = \mathcal O(1),$ meaning by the master
    theorem we have a time complexity of $\mathcal O(n).$

    \newpage

  \item Similarly, we can write the psuedocode as
    \begin{algorithm}
      \caption{D\&C Prefix Sum}
      \begin{algorithmic}
        \Procedure{PrefixSum}{$A[1\ldots n]$}
        \If{$n = 1$}
        \State \Return $A$
        \EndIf

        \vspace{5mm}

        \State $L \gets \text{PrefixSum}(A[1 \ldots n/2])$
        \State $R \gets \text{PrefixSum}(A[n/2 + 1 \ldots n])$

        \For{$i \in n/2 + 1 \ldots n$}
        \State $R[i] \gets R[i] + L[n/2]$
        \EndFor
        \EndProcedure
      \end{algorithmic}
    \end{algorithm}

    For $n = 1,$ the prefix sum of $A$ is $A$ itself (since there is
    one element), meaning our algorithm is correct for input size $1.$

    Assume our algorithm is correct for all input sizes $1 \leq k
    \leq n - 1.$ Then, consider input $A[1 \ldots n].$ By our
    inductive hypothesis,
    \begin{align*}
      \text{PrefixSum}(A[1 \ldots n/2])[j] &= \sum_{i = 1}^{j} A[i], \\
      \text{PrefixSum}(A[n/2 + 1 \ldots n])[j] &= \sum_{i = 1}^{j}
      A[n/2 + i], \\
    \end{align*}

    Then, our algorithm returns an array that satisfies
    \begin{align*}
      R[i] =
      \begin{cases}
        \sum_{i = 1}^{j} A[i] & 1 \leq i \leq n/2 \\
        \sum_{i =1}^{n/2} A[i] + \sum_{i = 1}^{j} A[n/2 + i] & n/2 +
        1 \leq i \leq n,
      \end{cases}
    \end{align*}
    which simplifies to
    \begin{align*}
      R[i] =
      \begin{cases}
        \sum_{i = 1}^{j} A[i] & 1 \leq i \leq n/2 \\
        \sum_{i = 1}^{j} A[i] & n/2 + 1 \leq i \leq n \\
      \end{cases}         = \sum_{i = 1}^{j} A[i],
    \end{align*}
    meaning this algorithm returns the prefix sum of $A$ for input size $n.$

    $\therefore,$ via strong induction we have proved that this
    algorithm is correct.

    The recurrence satisfies
    \begin{align*}
      T(n) = 2T(n/2) + \mathcal O(n), \qquad T(1) = \mathcal O(1),
    \end{align*}
    and via the Master theorem the asymptotic complexity is $\mathcal
    O(n \lg n).$

    \newpage

  \item This problem can simply be solved via a D\&C max algorithm.
    \begin{algorithm}
      \caption{D\&C Peak Finding}
      \begin{algorithmic}
        \Procedure{PeakFinding}{$A[1 \ldots n]$}
        \If{$n = 1$}
        \State \Return $A[1]$
        \EndIf

        \vspace{5mm}

        \If{$A[n/2] < A[n/2 + 1]$}
        \State \Return PeakFinding($A[n/2 + 1 \ldots n]$)
        \EndIf

        \vspace{5mm}

        \State \Return PeakFinding($A[1 \ldots n/2]$)
        \EndProcedure
      \end{algorithmic}
    \end{algorithm}
    To prove correctness, when $n = 1$ the peak is always $A[1],$
    meaning our algorithm is correct for input size $1.$

    Assume the algorithm is correct for all input sizes $1 \leq r
    \leq n -1.$ Let the peak be at index $k.$ Then, we have two cases:
    \begin{itemize}
      \item If $1 \leq k \leq n/2,$ then $A[i] > A[j]$ for all $j > i
        \geq k.$ Hence, since $n/2 + 1 > n / 2 \geq k,$ we can
        conclude that $A[n/2] > A[n/2 + 1],$ meaning we skip all
        branches and return $\text{PeakFinding}(A[1 \ldots n/2]).$ By
        our inductive hypothesis, $A[1 \ldots n/2]$ is in the form of
        a valid input with size $1 \leq n/2 \leq n - 1$, with a peak
        at index $k,$ meaning this will return the peak $A[k].$

      \item Otherwise, if $n/2 + 1 \leq k \leq n,$ then for all $i <
        j \leq k,$ we have $A[i] < A[j].$ As a result, we conclude
        that $A[n/2] < A[n/2 + 1],$ meaning we take the first branch
        and return PeakFinding($A[n/2 + 1\ldots n]$). Since $A[n/2 +
        1 \ldots n]$ is a valid input with size $1 \leq n/2 \leq n -
        1$ and peak at $k$, our inductive hypothesis means this will
        return $A[k].$
    \end{itemize}

    In either case, we return the correct peak $A[k],$ meaning our
    algorithm is correct for input size $n.$

    $\therefore$ by strong induction, our algorithm is correct for
    all input sizes $n \geq 1.$

    This satisfies recurrence $T(n) = T(n/2) + \mathcal O(1),$
    meaning we have an asymptotic complexity via the Master theorem
    of $\mathcal O(\lg n).$

    \newpage

  \item When our divide and conquer has the maximum sum subarray
    straddled between two subproblems, we would like to find the
    maximum suffix sum of the left subarray, and the maximum prefix
    sum of the right subarray. Hence, let us write an algorithm that
    computes the max subarray sum, the max prefix sum, and the max suffix sum.

    \begin{algorithm}
      \caption{Max Sum}
      \begin{algorithmic}
        \Procedure{m}{$A[1 \ldots n]$}
        \If{$n = 1$}
        \State \Return $A[1], A[1], A[1], A[1]$
        \EndIf

        \State $L, L_p, L_s, L_\sigma \gets \text{m}(A[1 \ldots n/2])$
        \State $R, R_p, R_s, R_\sigma \gets \text{m}(A[ n/2 + 1 \ldots n])$

        \vspace{5mm}

        \State $K \gets \max\{L, R, L_s + R_p\}$
        \State $P \gets \max\{L_p, L_\sigma + R_p\}$
        \State $S \gets \max\{R_s, R_\sigma + L_s\}$

        \vspace{5mm}

        \State \Return $K, P, S, L_\sigma + R_\sigma.$
        \EndProcedure
      \end{algorithmic}

      \begin{algorithmic}
        \Procedure{MaxSum}{$A[1 \ldots n]$}
        \State \Return $\text{m}(A[1 \ldots n])[0]$
        \EndProcedure
      \end{algorithmic}
    \end{algorithm}

    First, let us prove that m is a correct algorithm, i.e m returns
    the max subarray sum, the max prefix sum, the max suffix sum, and
    the total sum.

    For $n = 1,$ all quantities are $A[1],$ meaning m is correct for
    input size $1.$ Now, let us assume m is correct for some input
    size $1 \leq k \leq  n - 1.$ Then:
    \begin{itemize}
      \item The max subarray can either be a subarray of $A[1 \ldots
        n/2],$ a subarray of $A[n/2 + 1 \ldots n],$ or straddled in between.

        In the first two cases, by our inductive hypothesis, $L$ or
        $R$ will be the max subarray sum, and since this is the
        maximum subarray sum of $A,$ it will exceed $L_s + R_p,$
        meaning we will correctly return the maximum subarray sum.

        In the third case, the subarray with the maximum sum will be
        composed of some suffix of $A[1 \ldots n/2],$ and some prefix
        of $A[n/2 + 1 \ldots n].$ Let us assume that this suffix and
        prefix that compose the maximum subarray isn't maximal. Then,
        we can construct a new subarray of $A$ with the respective
        maximal suffixes and prefixes replacing the chosen suffix and
        prefix that has a larger sum, contradicting the maximality of
        the subarray.

        Thus, the maximal subarray must be composed of the maximal
        suffix of $A[1 \ldots n/2],$ and the maximal prefix of $A[n/2
        + 1 \ldots n].$ By our inductive hypothesis, $L_s$ is this
        maximal suffix, and $R_p$ is this maximal prefix, meaning we
        correctly return $K = L_s + R_p$ in this case.

        Thus, we know for input size $n,$ we correctly return the
        maximal subarray.

      \item The maximal prefix of $A[1 \ldots n]$ can either extend
        up till index $n/2$, or extend further than this. In the
        first case, since by our inductive hypothesis $L_p$ is this
        maximal prefix sum, $L_\sigma$ is the sum of the left array,
        and $L_p$ must exceed the other prefix $L_\sigma + R_p,$ we
        correctly return the maximal prefix sum.

        In the second case, there is no prefix of the left array that
        is maximal. Thus, the prefix must extend past index $n / 2,$
        meaning it must contain the entirety of the left array. Then,
        assume it does not consist of the maximal right prefix. Then,
        we can replace whatever prefix it does consist of with the
        maximal right prefix, yielding a larger prefix sum,
        contradicting it's maximality. Hence, the maximal prefix sum
        of $A[1 \ldots n]$ must consist of the entirety of the left
        array, and the maximal prefix sum of the right array. By our
        inductive hypothesis, $L_\sigma$ is the left array's sum, and
        $R_p$ is the maximal prefix of the right subarray. Since
        $L_\sigma + R_p$ is guaranteed to be at least as large as any
        prefix of the left array, we return the correct maximal
        prefix sum of $A[1 \ldots n]$ in the second case too.

        In both cases, input size $n$ returns the correct maximal
        prefix of $A[1 \ldots n].$

      \item WLOG, we can make the same argument for the maximal suffix sum.

      \item Since the sum of $A[1 \ldots n]$ is composed of the sum
        of the left array and the sum of the right array (which are
        $L_\sigma$ and $R_\sigma$ by our inductive hypothesis), for
        input size $n,$ m correctly returns the sum of $A[1 \ldots n].$
    \end{itemize}

    $\therefore$, by strong induction m correctly returns all four
    statistics. Thus, since MaxSum($A[1 \ldots n]$) returns m$(A[1
    \ldots n])[0],$ MaxSum is also correct.

    The recurrence of this algorithm satisfies $T(n) = 2T(n/2) +
    \mathcal O(1),$ meaning it's asymptotic complexity is $\mathcal O(n).$
\end{enumerate}
\end{document}
